\SourceLeaf{64}
possunt. Tunc igitur necessarium est, ut sit $n$ numerus fractus, isque verus, sin minus
series omnes, cum divergant, usu prorsus destituuntur.

9. Hucusque, in sola evolutione seriei $(1+n\cos\varphi)^{\nu}$ morati, integrationem formulae
propositae $\mathrm{d}\varphi\,(1+n\cos\varphi)^{\nu}$ negleximus. Illa autem, inventa serie ipsa,
nullo absolvitur negotio. Si enim fuerit:
\[
(1+n\cos\varphi)^{\nu} = A + A^{(1)}\cos\varphi + A^{(2)}\cos 2\varphi + .\,.\,.\,. + A^{(r)}\cos r\varphi + .\,.\,.\,.
\]
fiet
\[
\int \mathrm{d}\varphi\,(1+n\cos\varphi)^{\nu} = A\varphi + A^{(1)}\sin\varphi + \tfrac{1}{2}A^{(2)}\sin 2\varphi + .\,.\,. + \frac{1}{r}A^{(r)}\sin r\varphi + .\,.\,.\,.,
\]
cuius seriei coefficientes ex praecedentibus noti sunt.

10. Subiunxit in fine Cap. VI \textsc{Eulerus} integrationem formulae
\[
\mathrm{d}\varphi\,\log(1+n\cos\varphi),
\]
quam ut priorem ita perficit, ut peculiari artificio $\log(1+n\cos\varphi)$ in seriem secundum
cosinus angulorum multiplorum evolvat, antequam integrationem ipsam adgrediatur. Revera autem
haec series, ut casus specialis, iam in illa continetur, quam pro $(1+n\cos\varphi)^{\nu}$
invenimus. Opus est in hoc negotio nota propositione: esse
\[
\log(1+x) = \frac{(1+x)^{0}-1}{0},
\]
cui consentienter
\[
\log(1+n\cos\varphi) = \frac{(1+n\cos\varphi)^{0}-1}{0}.
\]
\par\noindent\hspace*{\parindent}In expressione igitur generali pro $(1+n\cos\varphi)^{\nu}$ coefficiens initialis unitate
mulctetur, dein in illo et ceteris omnibus ponatur $\nu=0$, et dividatur quivis per ipsam $0$.

$\alpha$) Est ipsius $A$, coefficientis initialis ad seriem $(1+n\cos\varphi)^{\nu}$ pertinentis, et
termino suo initiali, qui est $=1$, iamiam mulctati, generalis $h^{\mathrm{tus}}$
\[
= \frac{1}{2^{2h}} . {}^{2h}\mathfrak{B}^{(h)} . {}^{\nu}\mathfrak{B}^{(2h)}\,n^{2h}.
\]
Quodsi in illo ponatur $\nu=0$, fiet
\[
\frac{1}{2^{2h}} . \frac{0.(-1).(-2).\,.(-2h+1)}{1.2.\,.\,.\,. 2h} . {}^{2h}\mathfrak{B}^{(h)}\,n^{2h},
\]

